In this section, we derive a simple algorithm for computing additive error low-rank approximations to a psd matrix using the length-squared sampling procedure of Chapter~\ref{chapter:fast-sample}. Applying a result of Drineas, Kannan, and Mahoney~\cite{drineas_fast_2006-1}, we compute an approximate basis for the top $k$-dimensional left singular subspace of $\bA$ from $O(k/\epsilon^2)$ columns sampled by squared column norms. Then, we use a well-known leverage score sampling guarantee to approximately project $\bA$ onto this subspace. Given entry-wise access to $\bA\in\R^{n\times n}$, our algorithm outputs matrices $\bX, \bY\in\R^{n\times k}$ satisfying
\[\|\bA - \bX\bY^\top\|^2_F\leq \|\bA - \bA_k\|^2_F + \epsilon\|\bA\|^2_F\]
with constant probability, using $O(nk/\epsilon^2)$ entries of $\bA$ and $O(nk^2/\epsilon^4)$ arithmetic in expectation.

Our sample complexity incurs an additional $1/\epsilon$ factor in comparison to the best known algorithm for additive error low-rank approximation: Bakshi, Chepurko, and Woodruff~\cite{bakshi_robust_2020} achieve an additive error low-rank approximation using only $\tilde O(nk/\epsilon)$ entries, where $\tilde O$ hides polylogarithmic factors in $n$, $k$, and $\epsilon$. Their result is sample-optimal, matching a $\Omega(nk/\epsilon)$ query lower bound of Musco and Woodruff~\cite{musco_sublinear_2019} up to polylogarithmic factors. Nonetheless, the algorithm presented here admits a simple analysis. We pose as an open question whether it is enough to sample only $O(k/\epsilon)$ columns by squared column norms, thus leading to an $O(nk/\epsilon)$ algorithm. The comparison is summarized in Table~\ref{tab:additive-lra-comparison}. 

\begin{table}[H]
\centering
\small
\renewcommand{\arraystretch}{1.25}
\begin{adjustbox}{max width=\textwidth}
\begin{tabular}{@{}L{0.10\textwidth}L{0.23\textwidth}L{0.35\textwidth}C{0.2\textwidth}C{0.2\textwidth}@{}}
\toprule
Result
&
Assumptions
&
Guarantee
&
Query complexity
&
Arithmetic complexity
\\
\midrule
\cite{bakshi_robust_2020}
&
Psd matrix with entry-wise access.
&
\(\bX,\bY\in\R^{n\times k}\) such that
\[
    \|\bA-\bX\bY^\top\|_F^2
    \leq
    \|\bA-\bA_k\|_F^2
    +
    \epsilon\|\bA\|_F^2 .
\]
&
\[
    \widetilde O\!\left(
        \frac{nk}{\epsilon}
    \right)
\]
&
\[
    \widetilde O\!\left(
        n\left(\frac{k}{\epsilon}\right)^{\omega-1}
    \right)
\]
\\
\midrule
This work
&
Psd matrix with entry-wise access. 
&
\(\bX,\bY\in\R^{n\times k}\) such that
\[
    \|\bA-\bX\bY^\top\|_F^2
    \leq
    \|\bA-\bA_k\|_F^2
    +
    \epsilon\|\bA\|_F^2 .
\]
&
\[
    O\!\left(
        \frac{nk}{\epsilon^2}
    \right)
\]
in expectation
&
\[
    O\!\left(
        \frac{nk^2}{\epsilon^4}
    \right)
\]
in expectation
\\
\bottomrule
\end{tabular}
\end{adjustbox}
\caption{Comparison of additive-error low-rank approximation guarantees for
positive semidefinite matrices. Here \(\omega\) is the matrix multiplication
exponent, and \(\widetilde O\) hides polylogarithmic factors in the relevant
parameters.}
\label{tab:additive-lra-comparison}
\end{table}


We begin with the following theorem due to Drineas, Kannan, and Mahoney~\cite{drineas_fast_2006-1}, which shows that one can compute a good approximation to the top $k$-dimensional left singular subspace of $\bA$ from $O(k/\epsilon^2)$ columns sampled by squared column norms.

\begin{lemma}[\cite{drineas_fast_2006-1}, Theorem~4]
\label{thm:dkm06}
    Let $\bA\in\R^{m\times n}$ be a nonzero matrix and $1\leq k\leq c\leq n$. Let $p_1, \ldots, p_n$ satisfy
    \[p_i = \frac{\|\bA_{*, i}\|^2_2}{\|\bA\|^2_F}.\]
    Now, suppose $i_1, i_2, \ldots, i_c$ are $c$ independent samples from the squared column-norm distribution on $\bA$, and let $\bC\in\R^{m\times c}$ be the matrix whose columns are given by
    \[\bC_{*, j} := \frac{1}{\sqrt{cp_{i_j}}}\bA_{*, i_j}.\]
    Finally, let $\bH\in\R^{m\times k}$ be a matrix containing the top-$k$ left singular vectors of $\bC$ as columns. We have the following:
    \begin{enumerate}[(a)]
        \item Fix $\delta\in(0,1)$ and let $\eta=1+\sqrt{8\log(1/\delta)}$. If $c\geq 4k\eta^2/\epsilon^2$, then with probability at least $1-\delta$,
    \[\|\bA - \bH\bH^\top\bA\|^2_F\leq \|\bA - \bA_k\|^2_F + \epsilon \|\bA\|^2_F.\]
        \item If $c\geq 4k/\epsilon^2$, then 
        \[\E\|\bA - \bH\bH^\top\bA\|^2_F\leq \|\bA - \bA_k\|^2_F + \epsilon \|\bA\|^2_F.\]
        \item If $c\geq 4/\epsilon^2$, then
        \[\E\|\bA - \bH\bH^\top\bA\|^2_2\leq \|\bA - \bA_k\|^2_2 + \epsilon \|\bA\|^2_F.\]
    \end{enumerate}
\end{lemma}

With constant probability, Lemma~\ref{thm:dkm06} produces a matrix $\bH$ with orthonormal columns such that $\bH\bH^\top\bA$, the projection of $\bA$ onto the subspace spanned by $\bH$, yields an additive error approximation to $\bA$. To produce an explicit factorization, we need to regress the columns of $\bA$ onto the basis $\bH$, i.e., compute $\bH^\top\bA$, so that setting $\bX = \bH$ and $\bY^\top = \bH^\top \bA$ yields an additive error approximation $\bX\bY^\top$. However, computing $\bH^\top\bA$ requires us to read all of $\bA$, so instead we estimate this matrix by sampling
rows of \(\bA\) according to the leverage scores of \(\bH\). We state and prove the needed guarantee below for completeness. Closely related results appear in Lemma 8 of \cite{drineas_fast_2006} and Lemma 2.2.14 of \cite{musco_power_2018}.

\begin{lemma}
\label{lem:probabilistic-right-factor}
Let \(\bA\in\R^{n\times n}\), and let
\(\bH\in\R^{n\times k}\) have orthonormal columns. Let $q$ be the leverage score distribution on $\bH$, where
\[
    q_i = \frac{\ell_i}{k},\qquad \ell_i = \|\bH_{i,*}\|^2_2.
\]
Suppose \(i_1,\ldots, i_t\) are independent samples from \(q\), and construct the matrix $\widehat\bZ\in\R^{k\times n}$ given by
\[
    \widehat{\bZ}
    :=
    \frac{1}{t}\sum_{r=1}^t
    \frac{\bH_{i_r,*}^{\top}\bA_{i_r,*}}{q_{i_r}}
    .
\]
Then for every \(\delta\in(0,1)\), with probability at least \(1-\delta\),
\[
    \|\bA-\bH\widehat{\bZ}\|_F^2
    \leq
    \|\bA-\bH\bH^\top\bA\|_F^2
    +
    \frac{k}{\delta t}\|\bA\|_F^2.
\]
\end{lemma}

\begin{proof}
For a sample \(i\sim q\), define
\[
    \bR_j=\frac{\bH_{i,*}^{\top}\bA_{i,*}}{q_{i}}.
\]
Then
\[
    \E[\bR]
    =
    \sum_{i:\ell_i>0} q_i
    \frac{\bH_{i,*}^{\top}\bA_{i,*}}{q_i}
    =
    \bH^\top\bA.
\]
Since the samples are independent, we have
\(
    \E\|\widehat{\bZ}-\bH^\top\bA\|_F^2
    \leq
    \frac{1}{t}\E\|\bR\|_F^2.
\)
Moreover,
\[
\begin{aligned}
    \E\|\bR\|_F^2
    &=
    \sum_{i:\ell_i>0}
    q_i
    \left\|
        \frac{\bH_{i,*}^{\top}\bA_{i,*}}{q_i}
    \right\|_F^2  \\
    &=
    \sum_{i:\ell_i>0}
    \frac{\|\bH_{i,*}^{\top}\bA_{i,*}\|_F^2}{q_i}\\
    &=\sum_{i:\ell_i>0}\frac{\|\bH_{i,*}\|_2^2\|\bA_{i,*}\|_2^2}{q_i}\\
    &={k\cdot \|\bA_{i,*}\|_2^2} \qquad\text{using \(q_i=\ell_i/k\).}
\end{aligned}
\]
Thus,
\[
    \E\|\bR\|_F^2
    \leq
    k\sum_{i=1}^n \|\bA_{i,*}\|_2^2
    =
    k\|\bA\|_F^2
\qquad\text{and}\qquad
    \E\|\widehat{\bZ}-\bH^\top\bA\|_F^2
    \leq
    \frac{k}{t}\|\bA\|_F^2.
\]
By Markov's inequality,
\[
    \|\widehat{\bZ}-\bH^\top\bA\|_F^2
    \leq
    \frac{k}{\delta t}\|\bA\|_F^2\qquad\text{with probability at least \(1-\delta\).}\tag{$\ast$}
\]

Finally, let \(\bP=\bH\bH^\top\). Since \(\bP\) is the orthogonal projector
onto \(\operatorname{range}(\bH)\), we have the orthogonal decomposition
\[
    \bA-\bH\widehat{\bZ}
    =
    (\bI-\bP)\bA+\bH(\bH^\top\bA-\widehat{\bZ}).
\]
Thus,
\[
    \|\bA-\bH\widehat{\bZ}\|_F^2
    =
    \|\bA-\bH\bH^\top\bA\|_F^2
    +
    \|\bH(\widehat{\bZ}-\bH^\top\bA)\|_F^2.
\]
Since \(\bH\) has orthonormal columns,
\(
    \|\bH(\widehat{\bZ}-\bH^\top\bA)\|_F
    =
    \|\widehat{\bZ}-\bH^\top\bA\|_F.
\)
Recalling $(\ast)$ completes the proof.
\end{proof}

We are now ready to state the main algorithm of this section.

\begin{algorithm}[H]
\caption{Sublinear PSD Low-Rank Approximation}
\label{alg:sublinear-psd-low-rank}

\noindent
\textbf{Input:} psd matrix
\(\bA\in\R^{n\times n}\), target rank \(k\in[n]\), accuracy
parameter \(\epsilon\in (0,1)\).

\smallskip
Let
\(
    c=\lceil Ck/\epsilon^2\rceil
    % \qquad
    % t=\left\lceil \frac{2k}{\epsilon}\right\rceil.
\)
and $t = \lceil C'k/\epsilon\rceil$ for some absolute constants $C, C'$.

\smallskip
Run the following procedure:
\begin{enumerate}[
    label=\arabic*.,
    leftmargin=2em,
    itemsep=0.35em,
    topsep=0.45em,
    parsep=0pt
]
    \item Run Algorithm~\ref{alg:sampling-procedure} $c$ times on $\bA$, producing
    \(
        i_1,\ldots,i_c
    \)
    sampled by squared column norms.

    \item Read the full columns
    \(
        \bA_{*,i_1},\ldots,\bA_{*,i_c}
    \)
    and form the matrix \(\bB\in\R^{n\times c}\) whose \(r\)-th column is
    \[
        \bB_{*,r}
        =
        \frac{\bA_{*,i_r}}{\|\bA_{*,i_r}\|_2}.
    \]

    \item Let \(\bH\in\R^{n\times k}\) have orthonormal columns spanning a top-\(k\) left singular subspace of \(\bB\).

    \item Let $q$ be the leverage score distribution on $\bH$ with
    \(
        q_i=\frac{\|\bH_{i,*}\|_2^2}{k}
    \)
    Sample indices
    \(
        j_1,\ldots,j_t \overset{\mathrm{i.i.d.}}{\sim} q.
    \)
    \item Query the full rows
    \(
        \bA_{j_1,*},\ldots,\bA_{j_t,*}.
    \)
    and set
    \[
        \widehat{\bY}
        =
        \frac{1}{t}\sum_{r=1}^t
        \frac{\bA_{j_r,*}^{\top}\bH_{j_r,*}}{q_{j_r}}
        \in\R^{n\times k}.
    \]

    \item Return
    \(
        \bX=\bH,
        % \qquad
        \bY=\widehat{\bY}.
    \)
\end{enumerate}
\end{algorithm}


\begin{theorem}
\label{thm:sublinear-psd-low-rank}
Let \(\bA\in\R^{n\times n}\) be a nonzero psd matrix, let
\(1\leq k\leq n\), and let \(\epsilon\in(0,1)\). 
There are absolute constants \(C,C'>0\) such that
Algorithm~\ref{alg:sublinear-psd-low-rank}, run with
\[
    c=\left\lceil \frac{Ck}{\epsilon^2}\right\rceil,
    \qquad
    t=\left\lceil \frac{C'k}{\epsilon}\right\rceil,
\]
returns matrices \(\bX,\bY\in\R^{n\times k}\) satisfying
\[
    \|\bA-\bX\bY^\top\|_F^2
    \leq
    \|\bA-\bA_k\|_F^2+\epsilon\|\bA\|_F^2
\]
with probability at least \(2/3\). The algorithm queries $O(nk/\epsilon^2)$ entries of $\bA$ and uses $O(nk^2/\epsilon^4)$ arithmetic operations in expectation.
\end{theorem}

\begin{proof}
Let
\(
    \eta=1+\sqrt{8\log 6}
\)
and fix constants
\[
    C\geq 16\eta^2,
    \qquad
    C'\geq 12.
\]
We first prove the approximation guarantee.

Let \(i_1,\ldots,i_c\) be the column indices sampled in Step 1 of
Algorithm~\ref{alg:sublinear-psd-low-rank}. Let \(\widetilde{\bC}\) denote the
matrix of sampled columns appearing in Lemma~\ref{thm:dkm06}. We have
\begin{align*}
    \widetilde{\bC}_{*,r}
    &=
    \frac{1}{\sqrt{c p_{i_r}}}\bA_{*,i_r}\\
    &=\frac{\|\bA\|_F}{\sqrt c}
    \frac{\bA_{*,i_r}}{\|\bA_{*,i_r}\|_2}\qquad\text{since \(p_i=\|\bA_{*,i}\|_2^2/\|\bA\|_F^2\)}\\
    &=
    \frac{\|\bA\|_F}{\sqrt c}\bB_{*,r}.
\end{align*}
Thus
\(
    \widetilde{\bC}
    =
    \frac{\|\bA\|_F}{\sqrt c}\bB,
\)
is a scalar multiple of $\bB$, so \(\widetilde{\bC}\) and \(\bB\) have the same top \(k\) left singular
subspaces. In particular, the matrix \(\bH\) computed in Step 3 is also a valid
choice of \(\bH\) for Lemma~\ref{thm:dkm06}.

Apply Lemma~\ref{thm:dkm06}(a) with accuracy parameter \(\epsilon/2\) and
failure probability \(\delta=1/6\). Since
\[
    c
    \geq
    \frac{16k\eta^2}{\epsilon^2}
    =
    \frac{4k\eta^2}{(\epsilon/2)^2},
\]
we have, with probability at least \(5/6\),
\[
    \|\bA-\bH\bH^\top\bA\|_F^2
    \leq
    \|\bA-\bA_k\|_F^2
    +
    \frac{\epsilon}{2}\|\bA\|_F^2.
    \tag{1}
\]

Now condition on \(\bH\) produced by the algorithm. Let
\(
    \widehat{\bZ}=\widehat{\bY}^{\top}.
\)
Applying Lemma~\ref{lem:probabilistic-right-factor} with
\(\delta=1/6\), we obtain, with probability at least \(5/6\),
\begin{align*}
    \|\bA-\bH\widehat{\bZ}\|_F^2
    &\leq
    \|\bA-\bH\bH^\top\bA\|_F^2
    +
    \frac{6k}{t}\|\bA\|_F^2\\
    &\leq\|\bA-\bH\bH^\top\bA\|_F^2
    +
    \frac{\epsilon}{2}\|\bA\|_F^2.
    \qquad\text{since \(t\geq 12k/\epsilon\)}. \tag{2}
\end{align*}
By a union bound, \((1)\) and \((2)\) occur simultaneously with
probability at least \(2/3\), in which case we have
\[
\begin{aligned}
    \|\bA-\bX\bY^\top\|_F^2
    &=
    \|\bA-\bH\widehat{\bY}^{\top}\|_F^2  \\
    &=
    \|\bA-\bH\widehat{\bZ}\|_F^2 \\
    &\leq
    \|\bA-\bH\bH^\top\bA\|_F^2
    +
    \frac{\epsilon}{2}\|\bA\|_F^2 \\
    &\leq
    \|\bA-\bA_k\|_F^2
    +
    \epsilon\|\bA\|_F^2.
\end{aligned}
\]
This proves the claimed additive error guarantee.

It remains to bound the query complexity and time complexity. The \(c\) calls to
Algorithm~\ref{alg:sampling-procedure} use \(O(nc)\) entry queries in
expectation. Step 2 reads \(c\) full columns of \(\bA\), using \(nc\) further
entry queries. Step 5 reads \(t\) full rows of \(\bA\), using \(nt\) entry
queries. Therefore the total expected number of entry queries to $\bA$ is
\[
    O(nc+nt)
    =
    O\!\left(
        n\frac{k}{\epsilon^2}
        +
        n\frac{k}{\epsilon}
    \right)
    =
    O\!\left(\frac{nk}{\epsilon^2}\right),
\]
since \(\epsilon\in(0,1)\).

For the arithmetic bound, forming and rescaling \(\bB\in\R^{n\times c}\) costs \(O(nc)\)
operations after the sampled columns have been read. To compute a top \(k\)
left singular subspace for \(\bB\) we compute \(\bB^\top\bB\in\R^{c\times c}\) and its svd, which requires 
\(
    O(nc^2+c^3)
\)
arithmetic. For \(c\leq n\), this is \(O(nc^2)\), hence
\(
    O\left({nk^2}/{\epsilon^4}\right).
\)
Computing all leverage scores of \(\bH\) costs \(O(nk)\). Finally, forming \(\widehat{\bY}\) from the \(t\) sampled rows costs
\(O(ntk)\) arithmetic operations, since each sampled row contributes an
\(n\times k\) outer product. This is 
\(
    O\left({nk^2}/{\epsilon}\right),
\)
which is dominated by \(O(nk^2/\epsilon^4)\) for \(\epsilon\in(0,1)\). The
expected arithmetic cost of the sampling procedure contributes \(O(nc)\), which
is also dominated. Thus the total expected arithmetic complexity is
\[
    O\!\left(\frac{nk^2}{\epsilon^4}\right).
\]

% If \(c>n\), then \(k/\epsilon^2=\Omega(n)\), and one may instead read all of
% \(\bA\) and compute a rank-\(k\) truncated SVD directly. This uses
% \(O(n^2)\subseteq O(nk/\epsilon^2)\) entry queries and
% \(O(n^3)\subseteq O(nk^2/\epsilon^4)\) arithmetic, so the stated bounds still
% hold after this standard fallback. This completes the proof.
\end{proof}